\section{Comparison of Environmental Data Resources}
\label{sec:comparison}

\subsection{Scope and Counting Method}
\label{subsec:comparison-scope-method}

Group-07 and Group-08 both collect environmental resources about Bangladesh,
but their folder structures and preparation workflows differ. The comparison
therefore separates original source files, converted derivatives,
Entity--Relationship model artifacts, supporting evidence and final database
records instead of adding them into one total and treating the result as
quality.

Only measurements with a reproducible counting definition are included. File
size is not treated as database quality, and rows from arbitrary Excel or CSV
files are not treated as directly comparable final database cardinalities. A
converted copy is identified as a derivative rather than another unique
dataset. Where the supplied evidence does not establish an equivalent
measurement, the comparison states this instead of using zero.

The Group-07 register, current source tree, extraction registry and executable
SQLite database provide the Group-07 measurements. Group-08 measurements use
the public Google Drive folder metadata available during the review, the
supplied resource workbook, the latest ERD image and the supplied cleaning and
search-history evidence. The Drive counts describe the public folder state
visible during that inspection.

\subsection{Verified Repository Comparison}
\label{subsec:verified-repository-comparison}

Table~\ref{tab:verified-repository-comparison} compares like categories and
keeps Group-08's original and converted collections separate. The inspected
public Group-08 snapshot contains 66 files in \texttt{Datasets}, 47 CSV
derivatives, 5 root-level Google Sheets, 2 final or previous ER Diagram PNG
files, 18 individual-model JPEG files and 4 resource or support files. These
categories total 142 visible leaf files. The file-type check finds 57 XLSX
files, 1 XLS file, 14 native Google Sheets, 47 CSV files, 2 PNG files, 18 JPEG
files and 3 PDF files.

\begin{table}[H]
\centering
\caption{Verified repository and resource comparison}
\label{tab:verified-repository-comparison}
\small
\begin{tabular}{L{0.24\textwidth}L{0.34\textwidth}L{0.34\textwidth}}
\toprule
Measure & Group-07 & Group-08 \\
\midrule
Main organisation style & Source organisations separated into government,
non-government and formatted-output branches & Dataset topics, with a separate
CSV-converted collection \\
Current resource register & 139 entries under 18 named source labels & 35
entries under 6 named source labels \\
Older or candidate register & Not added to the current register count & 128
entries under 11 labels in a separate older sheet; not added to the current 35
\\
Original or current dataset files & 75 collected source or data files under the
stated source-corpus definition & 66 files in \texttt{Datasets}: 56 XLSX, 9
native Google Sheets and 1 XLS \\
Converted or extracted derivatives & 62 named extraction blocks from 9 loader
files; broader formatted outputs are kept separate from original sources & 47
CSV files in the separate converted collection \\
Source mapping and support evidence & Resource register, mapping workbooks,
exclusion records, source publications and quality logs & Resource workbook,
cleaning-problem PDF and two search or browsing-history PDFs \\
ER or model artifacts & Approved final ERD and schema definition file; no
repository-total artifact count is used & 2 ER Diagram PNG files and 18
individual-model JPEG files \\
Files reaching an implementation & 9 structured files produce 62 extracted
source blocks processed by the pipeline & Not established from the supplied
evidence \\
\bottomrule
\end{tabular}
\end{table}

Group-07 mainly keeps source ownership visible through organisation folders and
retains substantial source and reference evidence. Group-08 makes many subjects
easy to locate through topic names, and its separate CSV collection makes 47
converted tables easier to inspect directly. Group-08 also retains resource,
search-history and cleaning evidence, although it is arranged differently. Its
original and converted files require a clear rule for identifying the canonical
source and its derivative.

\subsection{Data Model and Implementation Comparison}
\label{subsec:model-implementation-comparison}

Table~\ref{tab:model-implementation-comparison} distinguishes a modelled entity
from an implemented database table. The latest supplied Group-08 ERD contains
29 modelled entity boxes by visual inspection. No executable Group-08 database
is present in the supplied evidence or public-folder snapshot. The comparison
therefore does not describe those entities as implemented tables and does not
invent database row, key or integrity results.

\begin{table}[H]
\centering
\caption{Data model and implementation comparison}
\label{tab:model-implementation-comparison}
\small
\begin{tabular}{L{0.23\textwidth}L{0.35\textwidth}L{0.34\textwidth}}
\toprule
Criterion & Group-07 & Group-08 \\
\midrule
Model evidence & Approved final ERD plus a schema definition file &
Latest integrated ERD plus 18 individual ER-model images \\
Executable implementation & Supplied SQLite database and build or verification
scripts & No executable database was found in the supplied or shared project
folder for equivalent database-level verification \\
Relation or entity count & 21 implemented tables and 2 views & 29 modelled
entities in the inspected ERD \\
Implemented schema size & 76 table columns and 28 logical foreign-key
relationships & Not established from the supplied evidence \\
Final database records and range & 730,324 rows covering 1948--2024 & Not
established from the supplied evidence \\
Database verification & \texttt{integrity\_check = ok}; 0 foreign-key
violations & Equivalent executable checks are not available \\
Source or data preparation & Source-specific extraction from 9 heterogeneous
files into 62 traceable blocks, followed by normalization and loading &
Topic-based original files, 47 CSV derivatives and documented cleaning examples
\\
\bottomrule
\end{tabular}
\end{table}

Group-07's source audit also reads 3,717,305 inspected source cells,
identifies 268,835 missing or unusable cells, approximately 7.2\%, and records
132,743 quality-problem occurrences across 50 classes. These values describe
its own extraction audit and are not presented as a Group-08 comparison because
no equivalent verified measurement is available.

Group-07's heterogeneous publications require more source-specific extraction
and cleaning, while Group-08's converted files provide convenient flat
representations for inspection. Group-07's supplied project proceeds further
into an executable relational implementation with automated integrity and
reproducibility checks. This observation applies only to the inspected
materials and does not claim that Group-08 has never built a database elsewhere.

\subsection{Limitations and Takeaways}
\label{subsec:comparison-limitations-takeaways}

The comparison is limited to the supplied project folders, attached evidence
and public Drive contents visible during inspection. Source-file counts do not
measure database quality. Original and converted copies are not separate unique
datasets unless that duplication is stated, and a modelled entity count is not
an implemented-table count.

Group-07's database-level measurements can be checked directly because an
executable SQLite database is supplied. Equivalent database-level statistics
are not assigned to Group-08 because no executable database was available in
the inspected material. Consequently, competency-query execution and timing
results are limited to Group-07's verified implementation.

The practical takeaway is that an environmental database project benefits from
clear source registration, preserved source evidence, structured or flat
derivatives where useful, explicit record grain, normalized relationships and
an implementation that another person can check independently. Only
measurements supported by the inspected evidence are compared here.
